%=====================%
%     COLOR NAMES     %
%=====================%

\definecolor{monokai-bg}{RGB}{250,250,250}     % Light gray background
\definecolor{monokai-comment}{RGB}{117,113,94} % Grayish comments
\definecolor{monokai-purple}{RGB}{137,89,168}  % Purple for keywords
\definecolor{monokai-green}{RGB}{58,110,59}    % Green for strings
\definecolor{monokai-orange}{RGB}{253,151,31}  % Orange for specials
\definecolor{monokai-blue}{RGB}{81,154,186}    % Blue for types/functions
\definecolor{graybg}{RGB}{249, 249, 249}
\definecolor{graybr}{RGB}{204, 204, 204}



%==============%
%     Math     %
%==============%

% --- Penambahan Teorema ---
\newtheorem{theorem}{\Teorema}              % Teorema
\newtheorem{lemma}[theorem]{\Lema}          % Lema
\newtheorem{definition}[theorem]{\Definisi} % Definisi
\newtheorem{example}[theorem]{\Contoh}      % Contoh

% --- Math Roman Alternatif ---
\newcommand{\upmath}[1]{\mathrm{#1}}

% --- Kumpulan Simbol Huruf Tegak ---
\newcommand{\deriv}{\upmath{d}}     % Turunan Diferensial
\newcommand{\adj}{\mathrm{adj}}     % Adjoin Matriks
\newcommand{\euler}{\mathrm{e}}     % Euler
\newcommand{\imaginary}{\mathrm{i}} % Imajiner

% Huruf Tegak Matematika a hingga z
\newcommand{\upa}{\mathrm{a}}
\newcommand{\upb}{\mathrm{b}}
\newcommand{\upc}{\mathrm{c}}
\newcommand{\upd}{\mathrm{d}}
\newcommand{\upe}{\mathrm{e}}
\newcommand{\upf}{\mathrm{f}}
\newcommand{\upg}{\mathrm{g}}
\newcommand{\uph}{\mathrm{h}}
\newcommand{\upi}{\mathrm{i}}
\newcommand{\upj}{\mathrm{j}}
\newcommand{\upk}{\mathrm{k}}
\newcommand{\upl}{\mathrm{l}}
\newcommand{\upm}{\mathrm{m}}
\newcommand{\upn}{\mathrm{n}}
\newcommand{\upo}{\mathrm{o}}
\newcommand{\upp}{\mathrm{p}}
\newcommand{\upq}{\mathrm{q}}
\newcommand{\upr}{\mathrm{r}}
\newcommand{\ups}{\mathrm{s}}
\newcommand{\upt}{\mathrm{t}}
\newcommand{\upu}{\mathrm{u}}
\newcommand{\upv}{\mathrm{v}}
\newcommand{\upw}{\mathrm{w}}
\newcommand{\upx}{\mathrm{x}}
\newcommand{\upy}{\mathrm{y}}
\newcommand{\upz}{\mathrm{z}}

% Huruf Tegak Matematika A hingga Z
\newcommand{\upA}{\mathrm{A}}
\newcommand{\upB}{\mathrm{B}}
\newcommand{\upC}{\mathrm{C}}
\newcommand{\upD}{\mathrm{D}}
\newcommand{\upE}{\mathrm{E}}
\newcommand{\upF}{\mathrm{F}}
\newcommand{\upG}{\mathrm{G}}
\newcommand{\upH}{\mathrm{H}}
\newcommand{\upI}{\mathrm{I}}
\newcommand{\upJ}{\mathrm{J}}
\newcommand{\upK}{\mathrm{K}}
\newcommand{\upL}{\mathrm{L}}
\newcommand{\upM}{\mathrm{M}}
\newcommand{\upN}{\mathrm{N}}
\newcommand{\upO}{\mathrm{O}}
\newcommand{\upP}{\mathrm{P}}
\newcommand{\upQ}{\mathrm{Q}}
\newcommand{\upR}{\mathrm{R}}
\newcommand{\upS}{\mathrm{S}}
\newcommand{\upT}{\mathrm{T}}
\newcommand{\upU}{\mathrm{U}}
\newcommand{\upV}{\mathrm{V}}
\newcommand{\upW}{\mathrm{W}}
\newcommand{\upX}{\mathrm{X}}
\newcommand{\upY}{\mathrm{Y}}
\newcommand{\upZ}{\mathrm{Z}}



%==============%
%     Code     %
%==============%

% --- Kode Sebaris Ber-highlight (Markdown Style) ---
\newcommand{\inlinesnippet}[1]{\fcolorbox{graybr}{graybg}{\lstinline‖#1‖}}
\newcommand{\verbsnippet}[1]{\fcolorbox{graybr}{graybg}{\verb|#1|}}



%===================%
%     Fixed Box     %
%===================%

\newcommand{\kotakMeterai}{
    \fbox{%
        \parbox[c][2.1cm][c]{2.9cm}{%
            \centering
            Meterai \\ Rp10.000
        }%
    }
}

%=======================%
%     More Commands     %
%=======================%

% --- Longkap Baris ---
\newcommand{\linesskip}[1]{\vspace{#1\baselineskip}}

% --- Keterangan Panjang (Lebih dari Sebaris)
\newcommand{\longcaption}[1]{\caption{\begin{tabular}[t]{@{}l@{}}#1\end{tabular}}}

% --- Sumber Lampiran Tabel ---
\newcommand{\tablesource}[1]{\vspace{.3\baselineskip}\caption*{Sumber: #1}\vspace{-\baselineskip}}
\newcommand{\tablesourceleft}[2]{
    
    \raggedright\medskip\hspace{#1}\small Sumber: #2}

% --- Sumber Lampiran Gambar ---
\newcommand{\figuresource}[1]{\vspace{-.9em}\caption*{Sumber: #1}\vspace{-.3\baselineskip}}

% --- Sumber Lampiran Kode ---
\newcommand{\lstsource}[1]{\begin{center}\vspace{-1.27\baselineskip}\setstretch{1}\small Sumber: #1 \vspace{.5\baselineskip}\end{center}}

% --- Link Email ---
\newcommand{\emailref}[1]{\href{mailto:#1}{#1}}

% --- Garis Titik-Titik ---
\newcommand{\dotsfill}[2]{
    \leavevmode%
    \makebox[#1]{%
        \cleaders\hbox to 0.44em{\smash{\lower#2\hbox{.}}\hss}\hfill%
    }%
}